\section{Exact supports and reduced product junctions}\label{st:syntax}
A support is a set of Newton indices with nonzero coefficients.
For the complete pure source these coefficients are nonnegative,
so supports can be computed without cancellations.

\begin{proposition}[Integer supports]\label{st:finite-support}
Every nonempty open complete support is an integer interval. If
$I_E=[l_E,u_E]\cap\ZZ$, then
\begin{align}
 I_h&=[0,M]\cap\ZZ,&
 I_{\T E}&=[d-u_E,d-l_E]\cap\ZZ,\notag\\
 I_{FG}&=
 \begin{cases}
 \varnothing,&l_F+l_G>d,\\
 [\max(l_F,l_G),\min(d,u_F+u_G)]\cap\ZZ,&l_F+l_G\le d.
 \end{cases}\label{st:interval}
\end{align}
Every coefficient in the displayed nonempty interval is positive.
An empty support remains empty under subsequent operations.
\end{proposition}
\begin{proof}
Only the product requires proof. For a proposed output $r$, the
allowable inputs lie in
$[l_F,\min(u_F,r)]$ and $[l_G,\min(u_G,r)]$.
Their sums fill an integer interval with lower endpoint $l_F+l_G$
and upper endpoint $\min(u_F,r)+\min(u_G,r)\ge r$.
Both $r$ and the lower endpoint are at most $d$, so this sum
interval meets $[r,d]$. Formula~\eqref{st:product} supplies a
positive summand. The same formula gives every necessary
inequality and the zero case.
\end{proof}
In particular
\[
 \phi(FG)=\sum_{r=0}^d [e_r]F\,[e_{d-r}]G.
\]
This identity also holds for arbitrary signed leaf coefficients.
Only the positive-support conclusion uses positivity.

\subsection{Strict real supports}
Scale the finite wire labels by $d$. A strict leaf label is
$0<a<1/2$. A strict retained product with inputs $x,y$ and output
$z$ satisfies
\begin{equation}\label{st:strict-gate}
 \max(x,y)<z<x+y<1.
\end{equation}
A turn sends $x$ to $1-x$. At a closed root impose $x+y=1$
instead of requiring a strict final product.

\begin{lemma}[The complete type table]\label{st:types}
The set of strictly realizable outputs of an open expression is
exactly one of
\[
 L=(0,1/2),\qquad H=(1/2,1),\qquad A=(0,1),\qquad\varnothing.
\]
The leaf has type $L$. Turn exchanges $L,H$, fixes $A$, and
preserves emptiness. Product is symmetric and has table
\[
 \begin{array}{c|ccc}
  &L&H&A\\\hline
 L&A&H&A\\
 H&H&\varnothing&H\\
 A&A&H&A
 \end{array}.
\]
The strict closed pairs are $AA,AL,AH,LH$ and their reversals.
\end{lemma}
\begin{proof}
For two low inputs and any $0<z<1$, choose both inputs in
$(z/2,\min(z,1/2))$. For an $L,A$ pair choose a sufficiently
small low input $c>0$ and the other input $z-c/2$.
Two $A$ inputs can use the first construction. A high input
forces $z>1/2$; for such $z$ choose that input $z-c/2$ and
the other input $c>0$ sufficiently small. Their sum is less
than one. Two high inputs cannot have sum less than one.
These choices prove sufficiency at every output in each stated
interval, as well as necessity. Finally, the closed condition is
$S_F\cap(1-S_G)\ne\varnothing$.
\end{proof}
\begin{definition}\label{st:strict-definition}
A cylindrical expression is \emph{strict} if it has a strict
reduced real state in this sense. This definition concerns its
original wire constraints.
\end{definition}
An empty strict type does not imply finite vanishing.
For example $\T h\,\T h$ has no strict state. At even $N$ its
support is empty, whereas at odd $N=2m+1$ its support
is the full interval $[m,d]$. Both original leaf
indices are then forced to their upper endpoint $m$,
so the contribution is boundary-only despite this
range of product outputs. Its cylindrical scalar
at output $d$ is positive. These assertions follow
from \eqref{st:interval}.

\subsection{A product with several inputs}
The symmetric tensor
$C_k(i_1,\ldots,i_k)=\phi(e_{i_1}\cdots e_{i_k})$
retains its normalization from \eqref{st:product}.
For $k=3$ its support is
\[
 i_1+i_2,\ i_1+i_3,\ i_2+i_3\le d,\qquad
 i_1+i_2+i_3\ge d.
\]
The input legs of a product gate are $x,y,1-z$ in this
symmetric notation.

\begin{lemma}[Multiple-input support]\label{st:hub-support}
For $n\ge2$ the product of basis elements with indices $i_1,\ldots,i_n$
is zero unless every pair satisfies $i_p+i_q\le d$. If all
these inequalities hold, its positive support is
\begin{equation}\label{st:many-support}
 [\max_p i_p,\min(d,\sum_p i_p)]\cap\ZZ.
\end{equation}
For fixed $0<x_p<1$, a prescribed output $z$ has a strict
resolution in any binary parenthesization precisely when
\begin{equation}\label{st:many-strict}
 x_p+x_q<1\ (p\ne q),\qquad
 \max_p x_p<z<\min(1,\sum_p x_p).
\end{equation}
A closed $k$-input product, $k\ge3$, has a strict resolution
precisely when
\begin{equation}\label{st:closed-hub}
 \sum_p x_p>1,\qquad x_p+x_q<1\quad(p\ne q).
\end{equation}
\end{lemma}
\begin{proof}
Associate a violating pair first to prove the finite zero
assertion. Otherwise split at any binary root. The lower
endpoints of the two child supports are their largest indices;
their sum is at most $d$ by the pair inequalities.
Proposition~\ref{st:finite-support} proves
\eqref{st:many-support} inductively, for this and every other
parenthesization.

For necessity in the strict statement, follow two inputs to
their first common gate. Its child labels are at least those
two input labels, and their sum is less than one. Following
outputs also gives the strict maximum and sum bounds.
For sufficiency, in a child set $S$ write
$L_S=\max_{p\in S}x_p$, $U_S=\min(1,\sum_{p\in S}x_p)$.
A nonsingleton child has output interval $(L_S,U_S)$;
a singleton has its prescribed value. Restrict the two
children to values below $z$. The lower endpoint of their
possible sum is $L_1+L_2<1$, and its upper endpoint is greater
than $z$, by \eqref{st:many-strict}. Thus their sum can be
chosen strictly between $z$ and one. If both are singletons,
their fixed sum already lies there. This gives a strict
parent gate, and induction fills its children.

For a closed product split at the removed final gate.
If both child sets have at least two inputs, their lower
endpoint sum is less than one and their upper endpoint sum
is greater than one, so their output intervals contain a
complementary pair. If one set is the singleton $x$, the
other interval contains $1-x$: its lower bound is smaller
by the pair inequalities, and its upper bound is larger
by the total-sum inequality. This proves sufficiency.
For necessity, two inputs meeting at a retained gate
have sum less than one. If their first common vertex
is the deleted closing gate, at least one branch
contains a strict product because $k\ge3$; its
output strictly exceeds the selected input. The two
selected inputs therefore still have sum less than
one. The same strict product makes the total of
the external inputs strictly larger than the two
closing outputs, whose sum is one.
For $k=2$ only that relation
remains, namely $x_1+x_2=1$.
\end{proof}

\subsection{The reduced junction tree}
Write $\mathbf J_{ij}=\mathbf1_{i+j=d}$. Raising one index of
$C_3$ by $\mathbf J$ recovers multiplication. Thus an internal
edge without a turn contracts through $\mathbf J$, while an
edge with a turn contracts through $\mathbf J\T=I$.
Contract unary chains modulo the algebraic identity $\T^2=I$,
suppress the degree-two closing vertex, and collapse every
internal edge without a turn. The result is a tree of
product vertices of valence at least three. We call these
vertices \emph{junctions}. Every remaining internal edge
has one turn. The original leaf labels are retained.
A leaf is \emph{flagged} if its incident edge has one turn:
it supplies $\T h$ rather than $h$.

\begin{theorem}[Reduced junction criterion]\label{st:reduced-junctions}
For a closed expression with $H\ge3$ leaves, the reduced
junction tree is unique under the indicated contractions.
The expression is strict if and only if each junction has
at most one flagged leaf. Its full finite source is the
contraction of the same tensors $C_k$, with all independent
leaf coefficients and repeated-root normal data retained.
\end{theorem}
\begin{proof}
The contracted components are exactly the connected components
of the internal edges without a turn, so they are uniquely
specified. Frobenius associativity identifies their tensor
with $C_k$. Lemma~\ref{st:hub-support} proves preservation of
strict feasibility at fixed external labels under every
collapse and resolution.

Two flagged leaves at a junction have labels greater than
$1/2$, contradicting \eqref{st:closed-hub}.
Conversely root at an internal edge, if one exists.
Every junction has at least two branches besides its parent.
An unflagged leaf supplies type $L$ and a flagged leaf type
$H$. Inductively an outward child product has type $H$ if
its own junction has one flagged leaf, and $A$ otherwise.
After the intervening turn its input to the parent has
type $L$ or $A$. The table of Lemma~\ref{st:types} therefore
gives exactly these claimed outgoing types. At the chosen
root edge one side is turned, leaving one of the closing
pairs $AA,AL,HA,HL$, each strict.
For a single junction, expose a leaf. With no flag the
remaining product has type $A$ and the pair is $L,A$.
With one flag expose it, giving $H,A$. The strict resolution
lemma transports the resulting state back to the literal
binary syntax.
\end{proof}
With two leaves there is no junction. The strict case has
one turn on their connecting edge and the two original leaf
labels are equal. We retain this case separately below.

\begin{figure}[ht]
\centering\begin{tikzpicture}[scale=.95,>=Latex,every node/.style={font=\small},
 junction/.style={circle,draw,fill=white,minimum size=6mm,inner sep=1pt}]
 \node[junction] (a) at (-1.65,0) {$C_3$};
 \node[junction] (b) at (1.65,0) {$C_3$};
 \coordinate (l1) at (-3.35,1.2);
 \coordinate (l2) at (-3.35,-1.2);
 \coordinate (l3) at (3.35,1.2);
 \coordinate (l4) at (3.35,-1.2);
 \draw[thick] (a)--(b) node[midway,above=3pt] {$\T$};
 \draw[thick] (a)--(l1) node[pos=.52,above left] {$\T$};
 \draw[thick] (a)--(l2);
 \draw[thick] (b)--(l3) node[pos=.52,above right] {$\T$};
 \draw[thick] (b)--(l4);
 \fill (l1) circle(2pt) node[left=4pt] {$h_1$};
 \fill (l2) circle(2pt) node[left=4pt] {$h_2$};
 \fill (l3) circle(2pt) node[right=4pt] {$h_3$};
 \fill (l4) circle(2pt) node[right=4pt] {$h_4$};
 \node at (0,-1.65)
 {$\phi\bigl((\T h_1)h_2\,\T((\T h_3)h_4)\bigr)$};
\end{tikzpicture}

\caption{A strict source with four original leaves. Each trivalent
vertex carries the symmetric product tensor. The middle edge has
a turn, so it is retained. A flagged leaf supplies $\T h_i$;
the other leaves supply $h_i$. Each vertex has one flagged
leaf. The labels $i$ remain independent when colors or crossings
are changed.}\label{st:tree-figure}
\end{figure}

For later estimates it is useful to record the exact outgoing
coupons. Root at a flagged leaf when one is present. A
nonroot junction gives
\begin{equation}\label{st:messages}
 Q_v=
 \begin{cases}
 \T\bigl(p_{i(v)}\prod_wQ_w\prod_lh_l\bigr),
       &\text{one flagged leaf},\\
 \T\bigl(\prod_wQ_w\prod_lh_l\bigr),
       &\text{no flagged leaf},
 \end{cases}
 \qquad p_i=\T h_i .
\end{equation}
The final scalar is $\phi(p_0\prod Q_w\prod h_l)$.
These are exact full-algebra identities, also for
$h_i=\T(R^{\eps_i}L_{j_i})$ at independent colors and signs.
The two outgoing types are respectively $L$ and $A$.
Commutativity and $\T^2=I$ in this calculation are not
claims that arbitrary physical turns or products can be
removed from the underlying tangle.

\begin{corollary}[Exposing any original leaf]\label{st:any-root}
Every original leaf of a strict closed tree with $H\ge3$
can be exposed as a closing character without changing
any other leaf. If its reduced edge parity is $\sigma$,
and $A$ is the product of the other branches at its
junction, then
\[
 Z_N=\phi((\T^\sigma h_j)A)
     =\phi(h_j\T^\sigma A)
     =\phi(p_j\T^{\sigma+1}A),\qquad p_j=\T h_j .
\]
The last factor is either an unturned product of at least
two $L/A$ branches, a decorated low module of the first
form in \eqref{st:messages}, or the turned output of an
unflagged product with at least two branches.
\end{corollary}
\begin{proof}
The pairing in \eqref{st:pairing} makes $\T$ self-adjoint:
$\phi((\T F)G)=\phi(F(\T G))$. This proves the identities.
If $\sigma=1$, the root leaf uses the only possible flag,
so all other branches are $L/A$ and
$\T^{\sigma+1}A=A$.
If $\sigma=0$ and another leaf at the junction is flagged,
$A=p_i\prod Q$ and $\T A$ is the decorated low module.
If there is no such flag, $\T A$ is the outgoing
unflagged product. The valence is at least three in
all cases. In particular an entirely unflagged
three-leaf tree is included.
\end{proof}
For a two-input product, two high Newton coordinates
vanish exactly: if $i,j>M$, then $i+j>d$ at both
parities. This finite observation is useful when the
last case of Corollary~\ref{st:any-root} has only
two inputs. It does not by itself bound a signed
character at the exposed root.
